\documentclass[11pt]{article}
\usepackage[T1]{fontenc}
\usepackage{lmodern,amsmath,amssymb,amsthm,mathtools}
\usepackage[margin=1in]{geometry}
\usepackage{microtype,enumitem,xcolor}
\usepackage[numbers,sort&compress]{natbib}
\usepackage{hyperref}
\hypersetup{colorlinks=true,linkcolor=blue!45!black,citecolor=blue!45!black,
urlcolor=blue!45!black,
pdftitle={The sharp Berry-Esseen bound for independent arrays with small Lyapunov ratio},
pdfauthor={Logan Bell}}
\setlength{\emergencystretch}{2em}
\numberwithin{equation}{section}
\newtheorem{theorem}{Theorem}[section]
\newtheorem{proposition}[theorem]{Proposition}
\newtheorem{lemma}[theorem]{Lemma}
\newtheorem{corollary}[theorem]{Corollary}
\theoremstyle{remark}
\newtheorem{remark}[theorem]{Remark}
\DeclareMathOperator{\Var}{Var}
\DeclareMathOperator{\Cov}{Cov}
\DeclareMathOperator{\supp}{supp}
\DeclareMathOperator{\Law}{Law}
\DeclareMathOperator{\Bern}{Bern}
\DeclareMathOperator{\sinc}{sinc}
\DeclareMathOperator{\dist}{dist}
\newcommand{\E}{\mathbb E}
\newcommand{\Pp}{\mathbb P}
\newcommand{\R}{\mathbb R}
\newcommand{\Z}{\mathbb Z}
\newcommand{\ind}{\mathbf 1}
\newcommand{\eps}{\varepsilon}
\newcommand{\be}{C_{\mathrm E}}
\newcommand{\ce}{c_{\mathrm E}}
\newcommand{\pe}{p_{\mathrm E}}
\newcommand{\qe}{q_{\mathrm E}}
\newcommand{\te}{\tau_{\mathrm E}}
\newcommand{\he}{h_{\mathrm E}}
\newcommand{\phz}{\phi(0)}
\title{The sharp Berry--Esseen bound for independent arrays\\
with small Lyapunov ratio}
\author{Logan Bell}
\date{13 September 2026\quad Version 0.5.0}
\begin{document}
\maketitle
\begin{abstract}
Every finite independent centered array with positive total variance
and finite third absolute moments satisfies $\Delta\le C_{\mathrm E}\ell$
whenever $\ell\le\exp(-\exp(20004))$, where $\Delta$ is the
Kolmogorov distance of the standardized sum from the normal law,
$\ell$ is its Lyapunov ratio, and
$C_{\mathrm E}=(3+\sqrt{10})/(6\sqrt{2\pi})$.
The coefficient is optimal. The proof develops the variational,
uniform-jitter, and local-comparison methods of He and Cheng for
summands with different laws. A finite approximate-period argument and
a quantitative form of Esseen's moment inequality identify a common
scale. Support variations put every coordinate of a maximizing array
into two narrow clusters or a small neighborhood of zero. An exact
comparison with a common-diameter Bernoulli array completes the proof.
All parameters are explicit. The common-diameter bound and an auxiliary
unrestricted bound are supplied by companion papers with validated
finite calculations.
\end{abstract}
\tableofcontents

\section{The theorem and the proof}\label{sec:theorem}
Let $X_1,\ldots,X_n$ be independent centered real random variables. Put
\begin{equation}\label{eq:basic-data}
 B^2=\sum_i\E X_i^2>0,\qquad T=\sum_i\E|X_i|^3<\infty,
 \qquad\ell=T/B^3,
\end{equation}
and write
\[
 \Delta=\sup_x\left|\Pp\left\{B^{-1}\sum_iX_i\le x\right\}-\Phi(x)\right|,
 \qquad\mathcal R(X_1,\ldots,X_n)=\Delta/\ell.
\]
Here $\Phi$ and $\phi$ are the standard normal distribution function
and density. Zero coordinates are allowed. Set
$\ce=3+\sqrt{10}$ and $\be=\ce\phz/6$.

\begin{theorem}\label{thm:main}
For every array satisfying \eqref{eq:basic-data},
\[
 \boxed{\quad\ell\le\ell_0:=\exp\bigl(-\exp(20004)\bigr)
          \quad\Longrightarrow\quad\Delta\le\be\ell.\quad}
\]
The coefficient $\be$ cannot be decreased, even when all coordinates
have the same two-point law.
\end{theorem}

The value of $\ell_0$ is a sufficient threshold. The proof does not
estimate the largest admissible threshold, and the sharp inequality
for unrestricted independent arrays remains open.

He and Cheng \cite[Theorem~1.1, Sections~2--5, and Appendix~A]{HeCheng2026} prove
the sharp inequality for iid sums once their sample size exceeds a
universal threshold. Their proof supplies three mechanisms used here:
coordinate contamination gives support equations, uniform jitter
reveals a lattice scale, and a local comparison separates small and
large residual variance. Independence with different laws requires
additional estimates. We prove these for the original coordinates,
including those outside the two-cluster class. In particular,
conditioning on a Bernoulli count does not make the residuals
independent; our conditional Fourier argument retains each label
together with its residual.

Two finite bounds from the companion papers enter the proof:
\begin{align}
 \mathcal R\bigl(h(B_1-p_1),\ldots,h(B_m-p_m)\bigr)&\le\be,
                                                        \label{eq:finite-input}\\
 \mathcal R(X_1,\ldots,X_n)&<9/20.                      \label{eq:upper-input}
\end{align}
In \eqref{eq:finite-input}, the $B_i$ are independent Bernoulli
variables with arbitrary parameters and positive total variance, and
$h>0$. This is the common-diameter theorem
\cite[Theorem~1.2(ii)]{BellFinite2026}; its proof uses Schulz's
binomial theorem \cite[Theorem~1]{Schulz2016} and validated finite
calculations. Equation~\eqref{eq:upper-input} follows from the
unrestricted bound \cite[Theorem~1.1]{BellUpper2026}.
These are the computational dependencies of the present theorem.
We also use Shevtsova's general-independent estimate
\begin{equation}\label{eq:published-excess}
 \Delta\le\be\ell+0.2538\ell^{4/3}\qquad(\ell\le10^{-3})
\end{equation}
\cite[Corollary~4.18, p.~303]{Shevtsova2012}. Its role is to localize
a maximizing counterexample. Its positive remainder is removed by
the subsequent support and local comparisons.

For the lattice step, define
\begin{equation}\label{eq:esseen-data}
 \pe=\frac{4-\sqrt{10}}2=\frac{1-1/\ce}{2},\quad
 \qe=1-\pe,\quad\te=\pe^2+\qe^2,\quad\he=1/\te.
\end{equation}
Esseen's moment inequality \cite[pp.~161--162]{Esseen1956} states
that a centered variable $Y$ supported on a translate of $h\Z$ satisfies
\begin{equation}\label{eq:esseen-moment}
 |\E Y^3|+3h\E Y^2\le\ce\E|Y|^3.
\end{equation}
For positive variance, equality holds precisely for
$Y\overset d=\pm h(B-\pe)$, $B\sim\Bern(\pe)$.
Appendix~\ref{app:confinement} proves the quantitative form needed here.
The iid asymptotic formula for this law gives a ratio tending to $\be$
and a Lyapunov ratio tending to zero
\cite{Esseen1956,MattnerShevtsova2019}. This proves optimality.

We use closed and strict CDFs $F_S(x)=\Pp(S\le x)$ and
$F_S^-(x)=\Pp(S<x)$. For a centered sum of variance $B^2$,
\begin{equation}\label{eq:branches}
 \sup_x|F_S(x)-\Phi(x/B)|
 =\sup_x\max\{F_S(x)-\Phi(x/B),\ \Phi(x/B)-F_S^-(x)\}.
\end{equation}
Reflection interchanges these branches. The comparisons below retain
these endpoint conventions and therefore include laws with atoms.

Section~\ref{sec:selection} selects an attained counterexample with
small Lyapunov ratio and bounds every support.
Section~\ref{sec:jitter} finds an approximate period by a finite
procedure. Section~\ref{sec:confinement} converts its moment information
into restrictions on every support point. Section~\ref{sec:local}
proves the independent local comparison with explicit parameters, and
Section~\ref{sec:completion} evaluates their composition. The appendices
give the Fourier estimates, the quantitative support argument, and the
conditional-moment estimates used in these sections.

\section{Selection and support equations}\label{sec:selection}
Normalize total variance to one, so $\ell=T$.
For a fixed number of coordinates, maximize the positive signed ratio
at a threshold, less a nondecreasing penalty on $\ell$.

\begin{lemma}[Localized selection]\label{lem:selection}
Let $0<\eta\le1/100$ and
$0<t\le\min\{1/2000,\eta^3/108\}$. If a finite array has
$\ell\le t$ and $\mathcal R>\be$, there is an attained penalized
maximizer with the same number of coordinates, a threshold $z$, and
\begin{equation}\label{eq:selected}
 \be<R<9/20,\quad\ell<2t,\quad
 R\le K:=R+\ell\psi_t'(\ell)\le R+\eta.
\end{equation}
Here $R$ is its full ratio. If a coordinate has variance $s$ and
third absolute moment $b$, then
\begin{equation}\label{eq:deletion-penalty}
 \psi_t\left(\frac{\ell-b}{(1-s)^{3/2}}\right)-\psi_t(\ell)
                         \le2\eta s.
\end{equation}
\end{lemma}

\begin{proof}
Put $f(v)=0$ for $v\le0$, $f(v)=v^2/(1+v)$ for $v>0$, and
$\psi_t(u)=2u^{1/3}f(\log(u/t))$. This penalty is continuously
differentiable, nonnegative, nondecreasing, and zero for $u\le t$.
After reflection if necessary, the supremum of
\begin{equation}\label{eq:objective}
              [F_S(z)-\Phi(z)]/\ell-\psi_t(\ell)
\end{equation}
over centered variance-one arrays of the specified length is above $\be$.
If the objective exceeds $r_0>0$, then $n^{-1/2}\le\ell\le r_0^{-1}$.
Bounded third moments give coordinatewise weak subsequences and uniform
integrability of squares, preserving means and total variance.
Chebyshev bounds the thresholds because their discrepancies are at
least $r_0/\sqrt n$. Moving-threshold Portmanteau for the closed CDF
gives the required upper semicontinuity at a limit. A strict loss of
third moments would increase the positive ratio and decrease the
penalty, contradicting the supremum. Thus the limit attains it.
This is the fixed-dimension argument of
\citet[Lemma~4.1]{HeCheng2026}, applied coordinatewise.

Since $f(\log2)>4/15>1/4$, one has $\psi_t(u)>u^{1/3}/2$ for
$u\ge2t$. For $u\le10^{-3}$ this exceeds the excess in
\eqref{eq:published-excess}; for $u\ge10^{-3}$ it is at least
$1/20>9/20-\be$. Hence the maximizer has $\ell<2t$.
For $t<u\le4t$, $f<1$ and $0\le f'\le1$ give
\[
 0\le u\psi_t'(u)\le\tfrac83(4t)^{1/3}
                         <3(4t)^{1/3}\le\eta.
\]
This holds also for $u\le t$. Each coordinate has
$s^{3/2}\le\ell<10^{-3}$, so $s\le1/4$ and
$\ell_-=(\ell-b)/(1-s)^{3/2}<(8/5)\ell<4t$.
If $\ell_->\ell$, integration on this actual interval gives
\[
 \psi_t(\ell_-)-\psi_t(\ell)
 \le\eta\log(\ell_-/\ell)
 \le\tfrac32\eta[-\log(1-s)]\le2\eta s.
\]
Otherwise monotonicity suffices. A larger discrepancy in the selected
array, reflected for the other branch, would increase its objective;
therefore $R$ is its full ratio.
\end{proof}

Write $s_i=\E X_i^2$, $b_i=\E|X_i|^3$,
$M_i=\E[X_i|X_i|]$, and let $G_i$ be the CDF of $S-X_i$.
\begin{lemma}[Support contact]\label{lem:contact}
For every $y\in\R$,
\begin{align}
 D_i(y):={}&G_i(z-y)-F_S(z)+\phi(z)y
       +\frac{z\phi(z)}2(y^2-s_i)\notag\\
 &+\frac32K\ell(y^2-s_i)
       -K\{|y|^3-b_i-3M_i y\}\le0.                 \label{eq:contact}
\end{align}
Equality holds at every $y\in\supp X_i$.
\end{lemma}
\begin{proof}
Replace the raw law of coordinate $i$ by
$(1-e)\Law(X_i)+e\delta_y$, retain the raw threshold, and center
and standardize the new array. The derivatives at $e=0$ of its mean,
variance, and total third absolute central moment are
$y$, $y^2-s_i$, and $|y|^3-b_i-3M_i y$.
The last differentiation is justified by a bound $C_y(1+X_i^2)$.
The raw CDF derivative is $G_i(z-y)-F_S(z)$. Differentiating
\eqref{eq:objective} and multiplying by $\ell$ gives
\eqref{eq:contact}, with its sign supplied by maximality.
Its integral under $\Law(X_i)$ is zero. Upper semicontinuity of $D_i$
then extends equality to the entire support. This is the
independent-coordinate form of the contamination argument of
\citet[Lemma~4.3]{HeCheng2026}.
\end{proof}

\begin{lemma}[Uniform support bound]\label{lem:support}
Every nonzero coordinate of the selected maximizer satisfies
\begin{equation}\label{eq:support}
 b_i<3\ell s_i,\qquad s_i<9\ell^2,
                 \qquad\supp X_i\subset(-8\ell,8\ell).
\end{equation}
\end{lemma}
\begin{proof}
Deletion, padded by a zero coordinate, stays in the maximization class.
Its ratio at every threshold is at most $R+2\eta s_i$.
Compare the convolution of the deleted normal approximation with $X_i$
to its convolution with $N(0,s_i)$. Put
$a=\phz(1+4\phz)/6<13/75$. Translation Taylor expansion and
$\E|N(0,s_i)|^3=4\phz s_i^{3/2}\le4\phz b_i$ give
\[
 R\ell(1-s_i)^{3/2}\le(R+2\eta s_i)(\ell-b_i)+ab_i.
\]
Using $1-(1-s_i)^{3/2}\le3s_i/2$ and rearranging gives
\[
 (R-a)b_i\le(3R/2+2\eta)\ell s_i.
\]
Since $R>2/5$, the right coefficient is less than $3(R-a)$.
This proves $b_i<3\ell s_i$, and Jensen gives $\sqrt{s_i}<3\ell$.

For $0\le s\le1/4$ put
\[
 H_s(z,y)=\Phi\left(\frac{z-y}{\sqrt{1-s}}\right)-\Phi(z)
                     +\phi(z)y+\tfrac12z\phi(z)(y^2-s).
\]
Translation and scale Taylor expansions give
\begin{align}
 |H_s(z,y)|&\le\phz|y|^3/6+\phz s|y|/2+s^2/6,    \label{eq:gaussian-global}\\
 |H_s(z,y)-\phi''(z)(sy/2-y^3/6)|
       &\le y^4/36+sy^2/6+s^2/6.                   \label{eq:gaussian-local}
\end{align}
Use $\|\phi''\|_\infty=\phz$, $\|\phi'''\|_\infty\le2/\pi<2/3$;
the scale remainder is bounded in absolute value by $4s^2/27<s^2/6$.
Substitute the deleted bound in the contact equality. Its constant
terms are
\[
 R(\ell_--\ell)-\tfrac32K\ell s+Kb
                  +[\psi_t(\ell_-)-\psi_t(\ell)]\ell_-.
\]
Since $(1-s)^{-3/2}<8/5$ and
$(1-s)^{-3/2}-1-3s/2\le6s^2$, every support point satisfies
\begin{equation}\label{eq:deletion}
 0\le H_s(z,y)-K|y|^3+\tfrac32K\ell y^2+3KM_i y
                       +6R\ell s^2+\eta b+4\eta\ell s.
\end{equation}
The contact equation first bounds the support, since its negative
cubic term dominates at infinity. Let $A_i$ be its maximum absolute
value. Use $|M_i|\le s$, $b\le sA_i$, $K\le23/50$, and
\eqref{eq:gaussian-global} to obtain
\[
 A_i^3/3\le(7/10)\ell A_i^2+(8/5)sA_i+(3/5)s^2+(1/25)\ell s.
\]
With $s\le9\ell^2$, $\ell\le1/8$, and $x=A_i/\ell$,
\[
 x^3/3\le7x^2/10+72x/5+1287/200.
\]
The difference is $2539/600>0$ at $x=8$ and has positive derivative
for $x\ge8$. Thus $A_i<8\ell$.
\end{proof}

\section{An approximate period and uniform jitter}\label{sec:jitter}
A uniform random shift suppresses the nonzero lattice peaks of a
characteristic function. Its period must be chosen from the array
itself. The following finite construction gives an explicit error.

\begin{lemma}[Finite period selection]\label{lem:period}
Let $\mu$ be a probability measure on $[-M,M]$, $M>0$, and put
\[
 g(u)=\int\frac{|e^{iud}-1|^2}{d^2}\,\mu(dd),\quad
 q(u)=\sqrt{g(u)},\quad\delta=1/(4M),
\]
where the integrand at zero is $u^2$. For $T\ge\delta$ and
$0<r\le\delta/4$ define
\[
 B_0=T/\delta+2,\quad J=\lceil\log_2(T/\delta)\rceil+1,\quad
 \eps=r/(2B_0^J),\quad E=\{u\in[\delta,T]:q(u)\le\eps\}.
\]
If $E\ne\varnothing$, there is $p\in[\delta,T]$ such that
$q(p)\le r/2$ and $\dist(u,p\Z)<r$ for every $u\in E$.
For $h=2\pi/p$ and $\ell>0$,
\begin{equation}\label{eq:high-frequency}
 \int_\delta^T e^{-g(u)/(2\ell^2)}|\sinc(hu/2)|\,\frac{du}{u}
 \le\frac{56\sqrt\pi}{9}\frac{r\ell}{p^2}
           +\log(T/\delta)e^{-\eps^2/(2\ell^2)}.
\end{equation}
If $E=\varnothing$, take $h=0$; the exponential term alone suffices.
\end{lemma}
\begin{proof}
The Hilbert-space triangle inequality gives subadditivity of $q$,
$q(-u)=q(u)$, and $q(u)\le|u|$. Thus $q$ is $1$-Lipschitz.
Also $g''\ge2-M^2g$ and $g(u)\ge u^2/2$ for $|u|\le\delta$.
Start with $p_0\in E$. If $v\in E$ has distance at least $r$ from
$p_j\Z$, choose the nearest integer $k$ and put
$p_{j+1}=|v-kp_j|$. While $p_j\ge\delta$,
\[
 r\le p_{j+1}\le p_j/2,\qquad
 q(p_{j+1})\le\eps+(T/\delta+1/2)q(p_j)\le B_0^{j+1}\eps.
\]
A proposed value in $[r,\delta]$ would satisfy both $q\le r/2$
and $q\ge r/\sqrt2$. There are at most
$\lceil\log_2(T/\delta)\rceil$ halvings before this happens.
The construction therefore stops with the required $p$.

Retain each $I_k=[kp-r,kp+r]\cap[\delta,T]$, $k\ge1$, that meets
$E$. On it $q\le\eps+2r\le5\delta/8$, so $g''\ge1$.
At its minimum $a_k$, including an endpoint minimum,
$g(u)\ge(u-a_k)^2/2$. Moreover
$|\sinc(\pi u/p)|\le r/u$ and $u\ge(k-1/4)p$.
Its integral is at most $2\sqrt\pi\ell r/[(k-1/4)^2p^2]$.
Sum using $\sum_{k\ge1}(k-1/4)^{-2}\le28/9$.
Outside the retained intervals $g>\eps^2$, proving the assertion.
\end{proof}

\begin{proposition}[Uniform jitter]\label{prop:jitter}
Suppose $\sum\E X_i^2=1$, $\ell=\sum\E|X_i|^3\le1/8$, and
$|X_i|\le8\ell$. For $0<\rho\le1$ put
\begin{align}
 \delta&=1/64,&T&=80/\rho,&r&=\rho\delta^2/48,\notag\\
 B_0&=T/\delta+2,&J&=\lceil\log_2(T/\delta)\rceil+1,
 &\eps&=r/(2B_0^J),\notag\\
 L_{\rm jit}(\rho)&=
 \min\left\{\frac18,\frac\rho{48000},
              \frac{\rho\eps^2}{8\log(T/\delta)}\right\}.
                                                        \label{eq:jitter-parameters}
\end{align}
If $\ell\le L_{\rm jit}(\rho)$, there is $0\le h\le128\pi$ such
that, for an independent $U_h$ uniform on $[-h/2,h/2]$,
\begin{equation}\label{eq:jitter}
 \left\|F_{S+\ell U_h}(x)-\Phi(x)
           -\frac\kappa6(1-x^2)\phi(x)\right\|_\infty\le\rho\ell,
 \qquad\kappa=\sum_i\E X_i^3.
\end{equation}
When $h>0$, $g(2\pi/h)\le r^2/4$ for the measure below. Here $U_0=0$.
\end{proposition}
\begin{proof}
Set $Y_i=X_i/\ell$, with independent copies $Y_i'$. The measure
\begin{equation}\label{eq:diff-measure}
 \mu(A)=\frac{\ell^2}{2}\sum_i
        \E[(Y_i-Y_i')^2\ind_{\{Y_i-Y_i'\in A\}}]
\end{equation}
is a probability measure on $[-16,16]$. It satisfies
\[
 g(u)=\ell^2\sum_i(1-|\E e^{iuY_i}|^2),\qquad
 |\E e^{iuS/\ell}|\le e^{-g(u)/(2\ell^2)}.
\]
Apply Lemma~\ref{lem:period}. The low-frequency and signed-smoothing
estimates in Appendix~\ref{app:fourier} bound the norm in
\eqref{eq:jitter} by
\begin{equation}\label{eq:jitter-error}
 12000\ell^2+12r\ell/\delta^2+
       \log(T/\delta)e^{-\eps^2/(2\ell^2)}+20\ell/T.
\end{equation}
Each term is at most $\rho\ell/4$ under
\eqref{eq:jitter-parameters}; for the third use $e^{-x}\le1/x$.
\end{proof}

\section{From moment stability to complete supports}\label{sec:confinement}
The contact equations strengthen the moment estimates into bounds for
each coordinate's full support.

\begin{proposition}[Two-cluster reduction]\label{prop:reduction}
Let $0<\lambda\le1$ and $N\ge1$. Set
\begin{equation}\label{eq:reduction-parameters}
 \rho=10^{-26}\lambda^2,\qquad\eta=10^{-10}\lambda,\qquad
 L_* =\min\{L_{\rm jit}(\rho),1/(2\sqrt N)\}.
\end{equation}
An attained maximizer of Lemma~\ref{lem:selection} with this penalty
bound and $\ell\le L_*$ can be represented, after division by $\ell$,
as independent coordinates
\begin{equation}\label{eq:clusters}
 Y_i=h_i(B_i-p_i)+U_i\quad(1\le i\le m),\qquad Z_1,\ldots,Z_s,
\end{equation}
where $m\ge N$, $h>0$, $B_i\sim\Bern(p_i)$, and
\begin{gather}
 p_i\in[1/4,3/4],\quad |h_i/h-1|\le\lambda,\quad
 \E(U_i\mid B_i)=0,\quad |U_i|\le\lambda h,              \label{eq:cluster-bounds}\\
 \E Z_j=0,\quad |Z_j|\le\lambda h,\quad
                        \sum_j\Var Z_j\le\lambda m h^2. \label{eq:extra-bounds}
\end{gather}
The representations are coordinatewise. Each pair $(B_i,U_i)$ is
independent of the other pairs and of all $Z_j$.
\end{proposition}

The proof is in Appendix~\ref{app:confinement}. Rounding at the period
from Proposition~\ref{prop:jitter} loses a controlled total second
moment. The deficit in Esseen's inequality forces that period close to
$\he$ and the maximizing threshold close to zero. Using the same
jitter for every deletion shows that its one-sided smoothing loss is
small at each support point. Comparing two such losses restricts every
difference of support points to a small neighborhood of $h\Z$.
The cubic expansion of \eqref{eq:deletion} bounds each scaled support
diameter by $7/2$ and puts the two-cluster probabilities in $[1/4,3/4]$.
Coordinates with small diameter have small total variance by moment
stability.
No lower bound on individual atom masses enters the argument.

\section{An exact comparison near a common diameter}\label{sec:local}
\begin{theorem}[Explicit independent local comparison]\label{thm:local}
Set
\begin{equation}\label{eq:local-constants}
 \mathcal A=e^{10000},\qquad\lambda=e^{-\mathcal A},
                   \qquad N=\lceil e^{\mathcal A}\rceil.
\end{equation}
Every independent array of the form
\eqref{eq:clusters}--\eqref{eq:extra-bounds}, with $m\ge N$, has
$\mathcal R\le\be$. Its ratio uses the total variance and the sum
of third absolute moments of all the original coordinates.
\end{theorem}

We follow the two residual regimes of
\citet[Section~3]{HeCheng2026}. The estimates below are uniform in
unequal probabilities, unequal gaps, and the additional coordinates.

\subsection{Normalization and conditional estimates}
Put $v_i=p_i(1-p_i)$ and define the variance-weighted mean gap
\[
 \bar h=\frac{\sum_i v_i h_i}{\sum_i v_i}.
\]
Then $(1-\lambda)h\le\bar h\le(1+\lambda)h$.
Divide every coordinate by $\bar h$, retaining $h_i,U_i,Z_j$ for the
rescaled gaps and residuals. Write
\begin{gather*}
 V=\sum_i v_i,\quad\mu=\sum_i p_i,
 \quad h_i=1+a_i,\quad\sum_i v_i a_i=0,\\
 K_0=\sum_iB_i,\quad A=\sum_i a_i(B_i-p_i),\quad
 U=\sum_iU_i+\sum_j Z_j,\quad W=A+U,\quad S=K_0+W,\\
 \Lambda_A=\sum_i v_i a_i^2,\quad
 \Lambda_U=\sum_i\E U_i^2+\sum_j\E Z_j^2,\quad
 \Lambda=\Lambda_A+\Lambda_U,\quad B^2=V+\Lambda.
\end{gather*}
The uncentered sum $S$ has mean $\mu$, and $\Cov(K_0,W)=0$.
Put $\eps_0=\max(\max_i\|U_i\|_\infty,\max_j\|Z_j\|_\infty)$,
with an empty maximum zero, and $\xi=\Lambda/V$.
The hypotheses imply
\begin{equation}\label{eq:local-domains}
 3m/16\le V\le m/4,\quad\max_i|a_i|\le3\lambda,\quad
 \eps_0\le2\lambda,\quad\xi\le8\lambda,\quad\Lambda\le2\lambda m.
\end{equation}
For example the upper bound before enlargement is
$\xi\le[\lambda+(7/4)\lambda^2]/[(3/16)(1-\lambda)^2]$.

Let $T_X$ and $\kappa$ be the total absolute and signed third central
moments of the original array. Its reference Bernoulli moments are
\[
 T_0=\sum_i v_i(p_i^2+(1-p_i)^2),\qquad
 \kappa_0=\sum_i v_i(1-2p_i).
\]
Conditional centering gives
\begin{align}
 |T_X-T_0|,\ |\kappa-\kappa_0|
      &\le2\sqrt{m\Lambda_A}+4\Lambda,                  \label{eq:moment-budget}\\
 T_X\ge V/3,\quad1/4\le T_X/B^2&\le2,\quad
 B^2\le m,\quad B/2\le P:=B^3/T_X\le4B.                \label{eq:prefactor}
\end{align}
Indeed $|(1+a)^3-1|\le4|a|$ and
$\sum v_i|a_i|\le\sqrt{V\Lambda_A}$.
Within each cluster the sign is fixed: its unperturbed magnitude is
at least $7/32>\eps_0$. Expanding the cubic changes either third
moment by at most $4\E U_i^2$, since the linear term vanishes.
An extra coordinate contributes at most $\eps_0\E Z_j^2$.
Jensen gives $T_X\ge(7/8)^3T_0\ge343V/1024>V/3$.
Every original coordinate has absolute value at most two, so
$T_X\le2B^2$, proving \eqref{eq:prefactor}.

For $m\ge10^{12}$ and $|k-\mu|\le6\sqrt m$,
Appendix~\ref{app:conditional} proves
\begin{align}
 |m_k|:=|\E(A\mid K_0=k)|&\le10^4\sqrt{\Lambda_A}/\sqrt m,
                                                        \label{eq:conditional-mean}\\
 |\Var(A\mid K_0=k)-\Lambda_A|&\le40000\Lambda_A/\sqrt m,\notag\\
 \E[(A-m_k)^4\mid K_0=k]&\le10^7\Lambda_A^2.             \label{eq:conditional-gap}
\end{align}
It also proves the relative label bound
\begin{equation}\label{eq:conditional-label}
 \left|\frac{\Pp(B_i=b\mid K_0=k)}{\Pp(B_i=b)}-1\right|
                         \le500/\sqrt m\qquad(b=0,1).
\end{equation}
For $\Lambda>0$ define
\begin{equation}\label{eq:DH}
 D=\Lambda^2+\eps_0^2\Lambda_U,\qquad
 H=\max\left\{\sqrt{\Lambda_A},
                   \frac\Lambda{\sqrt{\Lambda+\eps_0^2}}\right\}.
\end{equation}
Then, in the same count range,
\begin{gather}
 \E(W\mid K_0=k)=m_k,\quad
 \Lambda/2\le\Var(W\mid K_0=k)\le2\Lambda,\notag\\
 \E[(W-m_k)^4\mid K_0=k]\le10^8D,\quad
 \E(|W-m_k|\mid K_0=k)\ge H/30000,                       \label{eq:full-conditional}\\
 \E(W^4\mid K_0=k)\le10^9D,\qquad\E W^4\le400D.\notag
\end{gather}
To see the extension, condition first on all labels. The noise then
has variance $T_{\mathbf b}\le4\Lambda_U$ and fourth moment at most
$3T_{\mathbf b}^2+\eps_0^2T_{\mathbf b}$.
Its mean given $K_0$ is zero, and \eqref{eq:conditional-label}
compares its variance to $\Lambda_U$. Combine these facts with
\eqref{eq:conditional-gap} and $|x+y|^4\le8(|x|^4+|y|^4)$.
Interpolation $\E|Y|\ge(\E Y^2)^{3/2}/(\E Y^4)^{1/2}$ supplies
the second term in $H$. Conditional Jensen given the labels supplies
the first from $A-m_k$. In particular $|m_k|\le10^4H/\sqrt m$
uniformly, even when $\Lambda_A$ is much smaller than $\eps_0^2/m$.
Finally $\max a_i^2\le16\Lambda_A/3$ gives
$\E A^4\le9\Lambda_A^2$, and hence the unconditional bound.

\subsection{Uniform approximation and central thresholds}
Let $J$ be the CDF of $S$ plus an independent uniform variable on
$[-1/2,1/2]$. If $T\ge10$ and $\xi(T+1)^2\le1/56$, put
\begin{align}
 E_m(T)={}&\frac{401(1+\log T)}{\sqrt m}
       +(1+\log T)\xi(T+1)^2\notag\\
       &+60000(1+\log T)\sqrt m\,e^{-m/262144}+56/T.       \label{eq:local-error}
\end{align}
Appendix~\ref{app:fourier} proves, for $z=(x-\mu)/B$,
\begin{equation}\label{eq:local-jitter}
 \sqrt m\left\|J(x)-\Phi(z)
             -\frac\kappa{6B^3}(1-z^2)\phi(z)\right\|_\infty\le E_m(T).
\end{equation}
The expansion includes the third moments of the additional coordinates.

Fix
\[
 a_0=e^{1000},\quad\mathcal C=e^{1004},\quad\mathcal B=e^{1006},
 \quad q=e^{-\mathcal C},\quad T=e^{\mathcal B}.
\]
The inequalities $\mathcal A>100\mathcal B$,
$\mathcal B>7\mathcal C$, and $\mathcal C>50a_0$ imply, for every
$m\ge N$,
\begin{equation}\label{eq:local-numerical}
 \xi(T+1)^2<1/56,\quad E_m(T)<q/1000,\quad
 m^{-1/2}<q/1000,\quad200\lambda<q/1000.
\end{equation}
For a direct check, the four terms in \eqref{eq:local-error} are at most
\[
 802\mathcal B e^{-\mathcal A/2},\quad
 64\mathcal B e^{-\mathcal A+2\mathcal B},\quad
 240000\mathcal B(262144)^2e^{-3\mathcal A/2},\quad56e^{-\mathcal B}.
\]
Each is below $q/4000$. The third uses $e^{-x}\le2/x^2$, so this
verification places no upper bound on $m$.

Shift the jitter by $\pm1/2$. Either normalized signed discrepancy is
at most
\begin{equation}\label{eq:local-envelope}
 \phi(z)\left\{\frac{B^2}{2T_X}
      +\frac{|\kappa|}{6T_X}|1-z^2|\right\}+4E_m(T)+m^{-1/2}.
\end{equation}
The normalized Taylor error is at most $1/(4B)<m^{-1/2}$, using
$\|\phi'\|_\infty<1/4$,
$\|((1-z^2)\phi(z))'\|_\infty<5/4$, and \eqref{eq:prefactor}.
The main term is at most $(3+z^2)\phi(z)$, below $1/100$ for
$|z|\ge4$. Thus the full bound there is below $1/24$.

\subsection{Small residual variance}\label{subsec:small-residual}
Suppose $0<\Lambda\le e^{-1000}$. If $\Lambda=0$, the result is
exactly \eqref{eq:finite-input}. We use a one-sided loss estimate:
if $\E Y=0$, $\E Y^4\le1/160000$, $f(0)=0$, and
$3/4\le f'\le5/4$ on $[-7/10,7/10]$, then for $|u|\le3/5$,
\begin{align}
 \Pp(Y>u)+f(u)&\ge\tfrac38\E|Y|-176000\E Y^4,\notag\\
 \Pp(Y<u)-f(u)&\ge\tfrac38\E|Y|-176000\E Y^4.             \label{eq:one-sided}
\end{align}
This is the clipping argument of
\citet[Lemma~3.2 and equation~(83)]{HeCheng2026}. Appendix~\ref{app:conditional}
gives its proof with these constants and the stated threshold range.

The count estimate in Appendix~\ref{app:fourier} gives
\begin{equation}\label{eq:central-count}
 b_k:=\Pp(K_0=k)\ge e^{-100}/\sqrt m
                         \qquad(|k-\mu|\le6\sqrt m).
\end{equation}
Write a threshold with $|z|<4$ as $k+u$, $|u|\le1/2$, where $k$
is a nearest integer. Then $|k-\mu|\le5\sqrt m$.
For $f_k(u)=[\Phi((k-\mu+u)/B)-\Phi((k-\mu)/B)]/b_k$,
we have $3/4\le f_k'\le5/4$ on $[-1,1]$.
Indeed the relative count error is at most $18e^{100}/\sqrt m$,
and the logarithm of the Gaussian density ratio is at most
$100\xi+30/\sqrt m$; both are below $1/100$.

The exact count decompositions are
\begin{align}
 F_S(k+u)&=F_{K_0}(k)-b_k\Pp(W>u\mid K_0=k)+e^+_{k,u},\notag\\
 F_S^-(k+u)&=F_{K_0}(k-1)+b_k\Pp(W<u\mid K_0=k)+e^-_{k,u}.
                                                        \label{eq:layers}
\end{align}
Each remainder is bounded in absolute value by
\[
 \sum_{j\ne k}b_j\Pp\{|W|\ge|j-k|-1/2\mid K_0=j\}
                         \le2\cdot10^{11}D/\sqrt m.
\]
For $|j-k|\le\sqrt m$, use $b_j\le3/\sqrt m$,
\eqref{eq:full-conditional}, and
$\sum_{r\ne0}(|r|-1/2)^{-4}<40$.
For more distant counts use $\E W^4\le400D$ and a denominator
at least $m^2/16$.

Apply \eqref{eq:one-sided} to $Y=W-m_k$ conditionally on $K_0=k$,
using $f_k(v+m_k)-f_k(m_k)$ and $v=u-m_k$.
The bounds above give $|m_k|<1/20$ and the required fourth-moment
bound. Also $e^{-100}/5\le Pb_k\le5$.
The resulting normalized loss is at least $e^{-120}H$.
Its fourth-moment and leakage errors are below $10^{14}D$, and its
conditional-mean error is at most $62500\sqrt{\Lambda_A}/\sqrt m$.

Let $R_0=\mathcal R(B_1-p_1,\ldots,B_m-p_m)\le\be$.
Equation~\eqref{eq:moment-budget} gives
\[
 \left|\frac{B^3/T_X}{V^{3/2}/T_0}-1\right|
                  \le32\sqrt{\Lambda_A}/\sqrt m+96\Lambda/m.
\]
The normal variance change costs at most $\xi/8$ before multiplication
by $P$. Combining these errors with \eqref{eq:layers}, either central
normalized branch is at most
\begin{equation}\label{eq:small-comparison}
 R_0-e^{-120}H+e^{40}
           \left\{D+\frac{\sqrt{\Lambda_A}+\Lambda}{\sqrt m}\right\}.
\end{equation}
The full definition of $H$ gives
\[
 D/H\le(\Lambda+\eps_0^2)^{3/2},\quad
 \sqrt{\Lambda_A}/H\le1,\quad
 \Lambda/H\le\sqrt{\Lambda+\eps_0^2}.
\]
Since $\eps_0\le e^{-1000}$ and $m\ge e^{1000}$, the error is below
$e^{40}(3e^{-1500}+2e^{-500})H<e^{-120}H/2$.
Thus the central comparison is strict for every positive residual
variance. Finally $R_0\ge\phz/4>1/12$: apply the count estimate at
an integer nearest $\mu$ and take half its jump. The bound $1/24$
at distant thresholds is also below $R_0$.

\subsection{Positive residual variance and interval mass}
\label{subsec:large-residual}
Suppose $\Lambda\ge e^{-1000}$ and $|z|\le4$.
Choose all integers $|k-x|\le\max(1,\sqrt\Lambda)$. There are at
least $\sqrt\Lambda$ of them, all within $6\sqrt m$ of $\mu$.
Tilt the Bernoulli probabilities exponentially to make their count
mean $k$, and project the gap residual off the count, as in
Appendix~\ref{app:conditional}. On $K_0=k$ the residual is
$m_k^0+R_k$, where under the tilted measure $R_k$ is a sum of
independent centered pair residuals and additional coordinates, with
\begin{equation}\label{eq:tilted-residual}
 L:=\Var_\theta R_k\in[\Lambda/2,2\Lambda],\quad
 \Var_\theta K_0\ge m/10,\quad |m_k^0|\le\sqrt\Lambda,
 \quad\Cov_\theta(K_0,R_k)=0.
\end{equation}
Each residual summand is bounded by $8\lambda$.
The within-cluster noise variances change relatively by at most
$120/\sqrt m$ under the tilt; the extra coordinates do not change.
Thus these assertions include their actual variances.

Apply Lemmas~\ref{lem:conditional-fourier} and~\ref{lem:minorant}
to both quarter-intervals $(1/4,1/2)$ and $(-1/2,-1/4)$, with
\[
 L_0=\tfrac12e^{-1000},\quad C_0=4+2e^{500},\quad A_g=4e^{500},
 \quad\omega=16e^{A_g^2/6},\quad T_g=4\omega.
\]
The shift $y=x-k-m_k^0$ satisfies $|y|\le C_0\sqrt L$, and
$A_g\ge C_0+(7/16)/\sqrt{L_0}$.
The interval minorant has Gaussian expectation at least
$\phi(A_g)/(4\omega\sqrt L)$.
With $\beta=7/64$, the three terms in
\eqref{eq:conditional-fourier-error} are at most
\[
 128e^{-\mathcal A/2},\qquad128e^{500-\mathcal A},\qquad
                         10^7e^{4a_0-\mathcal A}.
\]
Their sum is below $e^{-10a_0}<\phi(A_g)/32$, because
$\phi(A_g)>e^{-9a_0}$, $\omega<e^{3a_0}$, and $T_g<e^{4a_0}$.
Also $8\lambda T_g<\beta/162$.
For the third term use $L\le m$, $e^{-x}\le2/x^2$, and $m\ge N$;
this is uniform over all possible residual variances.

Each conditional quarter-interval probability is at least
$\phi(A_g)/(8\omega\sqrt L)$. Multiply by the count probabilities
\eqref{eq:central-count} and sum over the chosen integers.
Since $L\le2\Lambda$ and $\mathcal C>50a_0$, this gives
\begin{equation}\label{eq:quarter-mass}
 \Pp\{S-x\in(1/4,1/2)\},\quad
 \Pp\{S-x\in(-1/2,-1/4)\}\ge q/\sqrt m.
\end{equation}
The two actual one-sided jitter losses therefore satisfy
\[
 J(x+1/2)-F_S(x)\ge q/(2\sqrt m),\qquad
 F_S^-(x)-J(x-1/2)\ge q/(2\sqrt m).
\]
After multiplication by $P\ge B/2$, each loss exceeds $q/10$.
Esseen's Bernoulli inequality gives $3V+|\kappa_0|\le\ce T_0$.
Using \eqref{eq:moment-budget}, the main term of
\eqref{eq:local-envelope} is consequently at most
\[
 \be+20\sqrt{\Lambda_A/m}+42\Lambda/m\le\be+200\lambda.
\]
Either central branch is below
$\be-q/10+200\lambda+4E_m(T)+m^{-1/2}<\be-q/20$.
The distant thresholds were already bounded by $1/24$.
This completes the proof of Theorem~\ref{thm:local}.


\section{Evaluation of the threshold}\label{sec:completion}
Take $\lambda,N$ from \eqref{eq:local-constants} and
$\rho,\eta,L_*$ from \eqref{eq:reduction-parameters}. A violation
with Lyapunov ratio at most
\begin{equation}\label{eq:tstar}
             t_*:=\min\{L_*/2,\eta^3/108,1/2000\}
\end{equation}
would give, by Lemma~\ref{lem:selection}, an attained violating
maximizer with $\ell<2t_*\le L_*$. Proposition~\ref{prop:reduction}
puts the entire array into the class of Theorem~\ref{thm:local},
contradicting $R>\be$.

It remains to show $\ell_0<t_*$. Write $A=\mathcal A=e^{10000}$,
so $\lambda=e^{-A}$, $N=\lceil e^A\rceil$,
$\rho=10^{-26}e^{-2A}$, and $\eta=10^{-10}e^{-A}$.
For the jitter parameters put
\[
 Q=T/\delta=5120/\rho,\quad B_0=Q+2,\quad
 J=\lceil\log_2Q\rceil+1,\quad r=\rho/(48\cdot4096),
 \quad\eps=r/(2B_0^J).
\]
Since $A>1000$, $\log2>2/3$, and $\log10<3$,
\[
 \log(1/\rho),\ \log Q,\ \log B_0,\ \log(1/r)<3A,
 \qquad J<5A.
\]
The ceiling is included: $J<\log Q/\log2+2<5A$. Thus
$\log(1/\eps)<3A+1+15A^2<16A^2$.
The negative logarithm of the third entry defining $L_{\rm jit}$ is
\[
 \log(1/\rho)+2\log(1/\eps)+\log(8\log Q)
                         <32A^2+4A<40A^2.
\]
The first two entries have smaller negative logarithms. Also
$N<2e^A$, so $\log(2\sqrt N)<A$ and $-\log L_*<40A^2$.
The other entries of \eqref{eq:tstar} have negative logarithms
\[
 3A+30\log10+\log108<4A,\qquad\log2000<A.
\]
Therefore $-\log t_*<41A^2$. Finally
$e^4>\sum_{j=0}^5 4^j/j!=643/15>41$, whence
\[
 -\log\ell_0=e^{20004}=e^4A^2>41A^2>-\log t_*.
\]
This completes the proof of Theorem~\ref{thm:main}.

\appendix
\section{Fourier estimates with numerical constants}\label{app:fourier}
We use $\widehat\nu(t)=\int e^{itx}\nu(dx)$ and
$\sinc u=\sin(u)/u$, with its continuous value at zero.
This appendix supplies the quantitative versions of the signed
smoothing and jitter arguments used in Sections~\ref{sec:jitter}
and~\ref{sec:local}.

\subsection{Signed smoothing and the low-frequency expansion}
If $F$ is a probability CDF with finite first absolute moment and $G$
is the CDF primitive of a finite signed measure of mass one and finite
first absolute moment, with Lipschitz constant $M_G$, then
\begin{equation}\label{eq:signed-smoothing}
 \|F-G\|_\infty\le\frac14\int_{-A}^A
       \frac{|\widehat F(t)-\widehat G(t)|}{|t|}\,dt+24M_G/A.
\end{equation}
Here $\widehat F$ and $\widehat G$ denote the transforms of their
measures. This is the signed smoothing lemma of
\citet[Lemma~2.1 and equation~(71)]{HeCheng2026}.
Its proof uses the even density
$k(x)=3\sinc(x/4)^4/(8\pi)$, whose transform is supported on
$[-1,1]$. The bound $k(x)\le3\min(1,256/x^4)/(8\pi)$ gives
$\E|Z|\le12/\pi<4$. With shift $a=16$, its one-sided tail
$\varepsilon_k$ is below $1/8$, and
$C_a=\E(16-Z)_+<18$.
At a positive near-maximum of $F-G$, monotonicity of $F$ gives
$(F-G)(x+y)\ge(F-G)(x)-M_Gy$ for $y\ge0$.
Average with the shifted kernel; use $-\|F-G\|_\infty$ on its
opposite tail. The negative branch uses the reflected shift.
This bounds the norm by
\[
 \frac{\|(F-G)*k_A\|_\infty}{1-2\varepsilon_k}
                    +\frac{C_aM_G}{A(1-2\varepsilon_k)}.
\]
Fourier inversion of the integrable CDF difference gives coefficients
$1/[2\pi(1-2\varepsilon_k)]\le1/4$ and
$C_a/(1-2\varepsilon_k)\le24$. Only $F$ needed monotonicity.
In particular the weaker coefficients $1/\pi$ and $40$ are valid.

Consider the variance-one row of Proposition~\ref{prop:jitter}.
Write $v_i=\E X_i^2$, $\kappa=\sum\E X_i^3$, and
$f_i(t)=\E e^{itX_i}$. For $|t|\le1/(64\ell)$,
\[
 |f_i(t)-1|\le v_it^2/2\le1/128,\quad
 \sum_i\E X_i^4\le8\ell^2,\quad\sum_i v_i^2\le64\ell^2.
\]
Using $|\log(1+z)-z|\le|z|^2$ for $|z|\le1/2$ yields
\[
 \sum_i\log f_i(t)=-t^2/2+\kappa(it)^3/6+w(t),\qquad
 |w(t)|\le(49/3)\ell^2t^4\le t^2/6.
\]
Thus the error against
$\widehat\nu(t)=e^{-t^2/2}(1+\kappa(it)^3/6)$ is at most
\begin{equation}\label{eq:log-error}
 \ell^2e^{-t^2/3}\{(49/3)t^4+t^6/72\}.
\end{equation}
Multiplication by $\sinc(h\ell t/2)$ adds at most
$(h^2\ell^2/24)(t^2+\ell|t|^5/6)e^{-t^2/2}$.
The positive-frequency integrals after division by $t$ are at most
$(591/8)\ell^2$ and $(13h^2/288)\ell^2$, respectively.
For $h\le128\pi$ their sum is below $7500\ell^2$.
The Gaussian comparison tail starting at $a=1/(64\ell)$ is at most
\[
 \ell^2\{64^2+64/3\}<4118\ell^2.
\]
Here use $\int_a^\infty e^{-t^2/2}dt/t\le a^{-2}$ and
$\int_a^\infty t^2e^{-t^2/2}dt\le2/a$.
The comparison CDF $\Phi(x)+\kappa(1-x^2)\phi(x)/6$ has
Lipschitz constant below $1/2$, since
$\|\phi'''\|_\infty\le2/\pi$ and $|\kappa|\le1/8$.
Applying the weaker smoothing coefficients $1/\pi,40$, and
Lemma~\ref{lem:period}, now proves \eqref{eq:jitter-error}.

For the local calculation a smaller low-frequency bound is convenient.
Put $m_{4,i}=\E X_i^4$ and
$A_i=-v_it^2/2+\E X_i^3(it)^3/6$.
On the same interval, $|A_i|\le(25/48)v_it^2<1/64$ and
\[
 |f_i(t)-e^{A_i}|\le m_{4,i}t^4/2.
\]
Symmetrization gives $|f_i(t)|\le e^{-v_it^2/3}$.
Telescope the products; omitting an individual decay factor costs
less than two. Since $\sum_i m_{4,i}\le8\ell^2$,
\[
 |\widehat F(t)-\widehat\nu(t)|
       \le8\ell^2t^4e^{-t^2/3}+(\ell^2t^6/72)e^{-t^2/2}.
\]
Integrating over both signs gives less than $73\ell^2$.
Using $|\sinc(h\ell t/2)-1|\le h^2\ell^2t^2/24$ and the actual
product bound adds at most $h^2\ell^2/8$. Hence
\begin{equation}\label{eq:small-low-frequency}
 \int_{|t|\le1/(64\ell)}
 \frac{|\sinc(h\ell t/2)\widehat F(t)-\widehat\nu(t)|}{|t|}\,dt
                         \le(73+h^2/8)\ell^2.
\end{equation}

\subsection{The count estimate}
\begin{lemma}\label{lem:count-estimate}
For a sum $K$ of independent Bernoulli variables, put
$\mu=\E K$ and $V=\Var K$. If $V\ge1/2$, then
\begin{equation}\label{eq:count-estimate}
 \left|\Pp(K=k)-\frac{\phi((k-\mu)/\sqrt V)}{\sqrt V}\right|
                   \le6/V,\qquad
 \sup_k\Pp(K=k)<1/\sqrt V.
\end{equation}
If the probabilities lie in $[\alpha,1-\alpha]$, let
$\beta=\alpha(1-\alpha)$ and $\gamma=\phi(M/\sqrt\beta)$.
For $m\ge144/(\beta\gamma^2)$ and $|k-\mu|\le M\sqrt m$,
$\Pp(K=k)\ge\gamma/\sqrt m$.
\end{lemma}
\begin{proof}
For $Z_i=B_i-p_i$ and $v_i=p_i(1-p_i)$, on $[-\pi,\pi]$,
\[
 |\E e^{itZ_i}|\le e^{-2v_i\sin^2(t/2)}
              \le e^{-2v_it^2/\pi^2},\qquad
 |\E e^{itZ_i}-e^{-v_it^2/2}|
              \le v_i(|t|^3/6+t^4/32).
\]
Telescope, retaining the decay from $V-v_i\ge V/2$.
The difference is at most
$V(1/6+\pi/32)|t|^3e^{-Vt^2/\pi^2}$.
Its integral divided by $2\pi$, plus the Gaussian tail, is at most
$(\pi^3/12+\pi^4/64+1/\pi^2)/V<6/V$.
Integrating the product modulus gives
$\sqrt\pi/(2\sqrt{2V})<1/\sqrt V$.
For the last assertion the error is at most half the central Gaussian
term, and $V\le m/4$ supplies the stated factor.
\end{proof}

\subsection{Conditioning an independent product on a count}
\begin{lemma}\label{lem:conditional-fourier}
Let $(B_i,R_i)$ be independent pairs, with
$B_i\sim\Bern(r_i)$, $r_i\in[a,1-a]$, $\E R_i=0$,
$|R_i|\le d_0$. Allow additional independent centered $Z_j$ with
$|Z_j|\le d_0$. Put
\[
 K=\sum_i B_i,\quad R=\sum_iR_i+\sum_jZ_j,\quad
 \E K=k\in\Z,\quad V=\Var K\ge1000,\quad L=\Var R>0,
\]
 and assume $\Cov(K,R)=0$. Set $\beta=a(1-a)$.
If $g$ is continuous and integrable, with integrable Fourier transform
supported on $[-T,T]$, and $d_0T\le\beta/162$, then
\begin{align}
 &\sup_y|\E[g(R-y)\mid K=k]-\E g(\sqrt L Z-y)|\notag\\
 &\qquad\le\frac{\|g\|_1}{\sqrt L}
 \left\{\frac{32}{\sqrt V}+\frac{8d_0}{\sqrt L}
                 +2T\sqrt{VL}e^{-\beta m/162}\right\},
 \qquad Z\sim N(0,1).                                  \label{eq:conditional-fourier-error}
\end{align}
\end{lemma}
\begin{proof}
Write $\chi(s,t)=\E e^{is(K-k)+itR}$.
For $|s|\le1/4$, $|t|\le T$, each centered individual phase $X$
is bounded by $1/2$. If $v=\E X^2$ and $b=\E|X|^3$,
symmetrization gives $|\E e^{iX}|\le e^{-v/3}$, and
$v^2\le b/2$ gives $|\E e^{iX}-e^{-v/2}|<b/4$.
The total phase variance is exactly $Vs^2+Lt^2$, by the zero covariance.
The third absolute moments sum to at most
$4(V|s|^3+d_0L|t|^3)$. Telescoping therefore gives
\[
 |\chi(s,t)-e^{-Vs^2/2-Lt^2/2}|
 \le2(V|s|^3+d_0L|t|^3)e^{-(Vs^2+Lt^2)/3}.
\]
Its double integral is at most
$56(VL)^{-1/2}(V^{-1/2}+d_0L^{-1/2})$;
use $\int e^{-u^2/3}du=\sqrt{3\pi}$,
$\int|u|^3e^{-u^2/3}du=9$.
For $1/4\le|s|\le\pi$, a Bernoulli factor is at most
$e^{-2\beta/81}\le1-\beta/81$ because $\sin(1/8)>1/9$.
The residual changes it by at most $Td_0\le\beta/162$, giving
$|\chi(s,t)|\le e^{-\beta m/162}$.
Its double integral is at most $4\pi T e^{-\beta m/162}$.
The omitted Gaussian count tail has integral at most
$8\sqrt{2\pi}(V\sqrt L)^{-1}e^{-V/32}$.

Fourier inversion in the count and residual uses
$|\widehat g|\le\|g\|_1$. By Lemma~\ref{lem:count-estimate},
$q=\Pp(K=k)$ differs from $q_0=(2\pi V)^{-1/2}$ by at most $6/V$;
$V\ge1000>288\pi$ ensures $q\ge q_0/2$.
Division by $q$ bounds the three preceding errors by
\[
 \frac{8\|g\|_1}{\sqrt L}(V^{-1/2}+d_0L^{-1/2}),\quad
 2\|g\|_1T\sqrt V e^{-\beta m/162},\quad
 \frac{3\|g\|_1}{\sqrt{VL}}e^{-V/32}.
\]
Replacing $q_0/q$ by one in the Gaussian expectation costs at most
$12\|g\|_1/\sqrt{VL}$. Enlarge $8+3+12$ to $32$.
The factors involving $y$ have modulus one, proving uniformity.
\end{proof}

\begin{lemma}[An interval minorant]\label{lem:minorant}
For an interval $(a,b)$ put $w_0=(a+b)/2$, $d=(b-a)/2$.
Fix $L_0>0$, $C_0\ge0$, and choose
\[
 A_g\ge C_0+(|w_0|+d/2)/\sqrt{L_0},\quad
 \omega=(2/d)e^{A_g^2/6},\quad T=4\omega.
\]
Then
\[
 g(w)=\left(1-\frac{(w-w_0)^2}{d^2}\right)
                      \sinc\bigl(\omega(w-w_0)\bigr)^4
\]
is at most $\ind_{(a,b)}$, has integrable Fourier transform supported
on $[-T,T]$, and satisfies
\[
 \|g\|_1\le4/\omega,\qquad
 \E g(\sqrt L Z-y)\ge\frac{\phi(A_g)}{4\omega\sqrt L}
            \quad(L\ge L_0,\ |y|\le C_0\sqrt L).
\]
If the bracket in \eqref{eq:conditional-fourier-error} is at most
$\phi(A_g)/32$, the corresponding conditional interval probability
is at least $\phi(A_g)/(8\omega\sqrt L)$.
\end{lemma}
\begin{proof}
On $|w-w_0|\le1/\omega$, the quadratic factor is at least $3/4$,
$\sinc\ge5/6$, and hence $g\ge1/4$. Its Gaussian density there
is at least $\phi(A_g)/\sqrt L$.
The negative part has integral at most $2/(d^3\omega^4)$, with
density at most $\phz/\sqrt L$. Since
$\omega^3=8\phz/(d^3\phi(A_g))$, subtraction leaves the asserted
lower bound. The integrals of
$\min(1,|\omega s|^{-4})$ and
$s^2\min(1,|\omega s|^{-4})$ are $8/(3\omega)$ and
$8/(3\omega^3)$; $\omega d\ge2$ proves the $L^1$ bound.
The transform of the fourth power of sinc is a compactly supported
cubic spline. Multiplication by the quadratic differentiates it twice,
leaving an integrable transform on the same support. The final assertion
follows from Lemma~\ref{lem:conditional-fourier}.
\end{proof}

\subsection{Jitter on the actual count lattice}
\citet[Lemma~A.4]{HeCheng2026} prove the iid jitter estimate.
We use the cancellation $\sum_i v_i a_i=0$ to control unequal gaps.
Use the normalization of Section~\ref{sec:local}, and let $f_i$ include
all original coordinate characteristic functions. Put
\[
 g(t)=\sum_i(1-|f_i(t)|^2),\qquad g_0(t)=2V(1-\cos t).
\]
For $\max|a_i|,\eps_0\le1/4$,
\begin{align}
 |g-g_0|&\le\Lambda t^2,\notag\\
 |g'-g_0'|&\le\Lambda(2|t|+2t^2),\notag\\
 |g''-g_0''|&\le\Lambda(2+8|t|+4t^2).                    \label{eq:resonance-derivatives}
\end{align}
To prove these bounds, couple each actual symmetrized difference with
$D_0=B_i-B_i'$, writing it as $D_0+E$, where
$E=a_iD_0+U_i-U_i'$. For an extra coordinate use $D_0=0$.
All interpolation points lie in $[-2,2]$, and
$\sum\E E^2=2\Lambda$.
Taylor-expand $1-\cos(td)$, $d\sin(td)$, and $d^2\cos(td)$ in $d$.
The terms involving $U$ vanish conditionally to first order, and the
remaining first-order terms cancel by $\sum_i v_i a_i=0$.
The second derivatives give exactly \eqref{eq:resonance-derivatives}.

If $T\ge10$ and $\xi(T+1)^2\le1/56$, each interval
$I_j=[2\pi j-1/4,2\pi j+1/4]$ meeting $[-T,T]$ has a unique
minimum $t_j$ of $g$, with
\[
 |t_j-2\pi j|\le4\Lambda(T+1)^2/V,\qquad
 \left|\prod_i f_i(t)\right|\le e^{-V(t-t_j)^2/4}\quad(t\in I_j).
\]
Indeed $g''\ge27V/16\ge V$ there, and the endpoint derivatives
have opposite signs: their reference magnitudes exceed $V/3$ while
the errors are at most $V/14$.
Outside these intervals $g\ge V/25$, giving product modulus
at most $e^{-V/50}$.
For $j\ne0$, the raw jitter multiplier $\sinc(t/2)$ vanishes at
$2\pi j$ and has Lipschitz constant at most $1/4$.
Using $|t|\ge5|j|$, integrate its linear bound against the displayed
Gaussian and sum the harmonic series. The result over both signs is
\begin{equation}\label{eq:resonance-integral}
 \sum_{j\ne0}\int_{I_j\cap[-T,T]}
       \frac{|\sinc(t/2)\prod_i f_i(t)|}{|t|}\,dt
 \le(1+\log T)\left\{\frac2{5V}
        +\frac{4\sqrt\pi}{5}\frac{\Lambda(T+1)^2}{V\sqrt V}\right\}.
\end{equation}
For $j\ge1$ the summand is at most
$1/(5jV)+\sqrt\pi|t_j-2\pi j|/(10j\sqrt V)$.

We now verify \eqref{eq:local-jitter}. Standardize by $B$ and put
$d=T_X/B^2\in[1/4,2]$. The Lyapunov ratio is $d/B$, and the
support bounds satisfy the hypotheses of
\eqref{eq:small-low-frequency}. Raw frequencies $|u|\le1/128$
lie in that estimate's range; raw jitter width one corresponds to
$h=1/d\le4$. Its integral is at most $300/B^2$, and after
\eqref{eq:signed-smoothing} and multiplication by $\sqrt m$ it
costs at most $400/\sqrt m$.
Equation~\eqref{eq:resonance-integral} costs at most
$(1+\log T)\{m^{-1/2}+\xi(T+1)^2\}$ in this normalization.
Between raw frequencies $1/128$ and $1/4$, the product is bounded
by $e^{-B^2u^2/3}$. Outside the resonance intervals it is bounded by
$e^{-V/50}$. The Gaussian comparison tail beginning at $B/128$ is
at most
\[
             (32854/B^2+1/192)e^{-B^2/32768}.
\]
Their total contribution after smoothing and scaling is below
$60000(1+\log T)\sqrt m e^{-m/262144}$, using $V\ge3m/16$.
Finally the cutoff $BT$ costs at most $24\sqrt m/(BT)<56/T$;
the standardized signed comparison has Lipschitz constant below one.
These are exactly the terms of \eqref{eq:local-error}.

\section{Quantitative confinement of the original supports}
\label{app:confinement}
This appendix proves Proposition~\ref{prop:reduction}. Its first two
lemmas turn a small lattice-moment deficit and a small smoothing loss
into explicit distances. The remaining argument applies them to the
contact equations and records the constants needed for the theorem.

\subsection{The deficit in Esseen's moment inequality}
We use a variance-weighted decomposition to quantify the deficit;
compare \citet[Lemma~A.9]{HeCheng2026} for stability of a single lattice law.
Put $c=\ce$, $p=\pe$, and $\tau=\te$ within this subsection.
Let $Y$ be centered, nondegenerate, and supported on a translate of
$h\Z$. Its oriented deficit is
\[
 \delta_Y=\frac{c\E|Y|^3-\E Y^3-3h\E Y^2}{h\E Y^2}.
\]
Write $P_-$ for the law of $-Y$ on $(0,\infty)$, $P_+$ for the
law of $Y$ there, and $m=\E Y_+=\E Y_->0$.
The centered two-point law $P_{u,v}$ at $-u,v$ has masses
$v/(u+v),u/(u+v)$. Direct integration gives
\[
 \Law(Y)=\Pp(Y=0)\delta_0+\int P_{u,v}\,\lambda(du,dv),
 \qquad\lambda(du,dv)=\frac{u+v}{m}P_-(du)P_+(dv).
\]
Its variance-weighted mixing measure
$\omega(du,dv)=uv\lambda(du,dv)/\E Y^2$ is a probability measure.
Set $n=(u+v)/h\in\{1,2,\ldots\}$ and $a=u/(u+v)$.
The identity
\[
 c\{a^2+(1-a)^2\}-(1-2a)=3+2c(a-p)^2
\]
gives the exact formula
\begin{equation}\label{eq:lattice-deficit}
 \delta_Y=\int\{3(n-1)+2cn(a-p)^2\}\,d\omega.
\end{equation}
This proves the oriented inequality and, by reflection, Esseen's
inequality \eqref{eq:esseen-moment}, including its equality case.

If $\delta_Y<3$, some component has $n=1$. Its two points are the
unique nearest lattice points straddling zero, say $-ah,(1-a)h$;
in particular this translate does not contain zero. Its variance
weight is at least $1-\delta_Y/3$. Every other component has variance
at least $h^2a(1-a)$, so this variance weight is at most its ordinary
mixture mass. Thus, for a centered law $Q$,
\begin{gather}
 \Law(Y)=(1-r_Y)\Law(h(B_a-a))+r_YQ,\quad r_Y\le\delta_Y/3,\notag\\
 |a-p|\le\sqrt{\frac{\delta_Y}{2c(1-\delta_Y/3)}},\quad
 h^2a(1-a)\le\E Y^2\le\frac{h^2a(1-a)}{1-\delta_Y/3}.   \label{eq:nearest-law}
\end{gather}
If $|Y|\le M$, coupling the mixture and then the two Bernoulli laws
bounds their Wasserstein distance by
\[
 W_1(\Law(Y),\Law(h(B_p-p)))
 \le\frac{2M}{3}\delta_Y
       +2h\sqrt{\frac{\delta_Y}{2c(1-\delta_Y/3)}}.
\]
Here $W_1$ is the infimum of $\E|Y-Z|$ over couplings.

For a collection of centered $h$-lattice variables, use one common
orientation and set
\[
 V=\sum_i\E Y_i^2,\quad B=\sum_i\E|Y_i|^3,\quad
 K=\sum_i\E Y_i^3,\quad \Delta_L=(cB-K-3hV)/(hV).
\]
Degenerate coordinates have zero variance weight. Applying the preceding
bound on $\delta_{Y_i}\le1$, the crude bound $M+h$ on its complement,
and Cauchy--Schwarz gives
\begin{equation}\label{eq:lattice-W1}
 \frac1V\sum_i\E Y_i^2 W_1(\Law(Y_i),\Law(h(B_p-p)))
 \le(5M/3+h)\Delta_L+h\sqrt{3/c}\sqrt{\Delta_L}.
\end{equation}
There are also useful moment-ratio estimates. Put
\[
 a_D=\sqrt{(1+D/3)D/(2c)},\quad
 b_D=\tau D/3+D/c+(2/c)a_D,\quad k_D=D/(3c)+2a_D.
\]
If $\Delta_L\le D$, then
\begin{equation}\label{eq:lattice-ratios}
 \left|\frac B{hV}-\tau\right|\le b_D,\qquad
 \left|\frac K{hV}-\frac1c\right|\le k_D,\qquad \frac B{hV}\ge\frac12.
\end{equation}
Indeed \eqref{eq:lattice-deficit} gives
$\E_\omega(n-1)\le D/3$,
$\E_\omega[n(a-p)^2]\le D/(2c)$, and
$|\E_\omega[n(a-p)]|\le a_D$.
Expand $n\{a^2+(1-a)^2\}$ and $n(1-2a)$ about $p$.
The last inequality follows directly from $n\ge1$.
All these formulas allow arbitrary variance weights, so they also hold
when each sum is multiplied by a common positive factor.

\subsection{Rounding at an approximate period}
For the selected array put $Y_i=X_i/\ell$.
Then $|Y_i|\le8$, $\E Y_i^2\le9$, and
\[
 \ell^2\sum_i\E Y_i^2=\ell^2\sum_i\E|Y_i|^3=1.
\]
At a nonzero frequency $u$, take the phase of
$\E e^{iuY_i}$, using any phase if the expectation vanishes.
Round $Y_i$ to the nearest point of the corresponding lattice of span
$h=2\pi/|u|$ and then center the rounded variable, calling it $Q_i$.
The inequality $|v|\le(\pi/2)|e^{iv}-1|$ for $|v|\le\pi$ gives
\begin{equation}\label{eq:rounding}
 e:=\ell^2\sum_i\E|Y_i-Q_i|^2\le h^2g(u)/8,\qquad |Q_i|\le8+h.
\end{equation}
In detail the uncentered squared rounding error is at most
$(h^2/16)\E|e^{iuY_i}-e^{i\theta_i}|^2
=(h^2/8)(1-|\E e^{iuY_i}|)$, which is no larger than the expression
using $1-|\E e^{iuY_i}|^2$. Centering decreases that error.
If $M=8+h$, Cauchy--Schwarz yields
\begin{align}
 \ell^2\sum_i|\E Q_i^2-\E Y_i^2|&\le2\sqrt e+e,\notag\\
 \ell^2\sum_i|\E|Q_i|^3-\E|Y_i|^3|,\quad
 \ell^2\sum_i|\E Q_i^3-\E Y_i^3|
                &\le(3M/2)\sqrt e(2+\sqrt e).           \label{eq:rounding-moments}
\end{align}
For the cubic bounds use
$||x|^3-|y|^3|,|x^3-y^3|\le(3M/2)|x-y|(|x|+|y|)$.
These bounds depend on total variance, not on the number of coordinates.

\subsection{Comparing two small smoothing losses}
\begin{lemma}\label{lem:support-spacing}
Let $F$ be the CDF of a finite positive measure and let
$J=F*\operatorname{Unif}[-h/2,h/2]$, $h>0$.
Suppose $J(x)$ differs from an affine function of slope $q>0$ by at
most $\varepsilon$ on $[-D_0-2h,3h]$.
Put $H(a)=J(a+h/2)-F(a)$. If $0\le d\le D_0$ and
$H(0),H(-d)\le s$, then, for $0<r_0\le h/4$ and
$K_0=\lceil D_0/h\rceil$,
\begin{equation}\label{eq:spacing-bound}
 \dist(d,h\Z)\le3r_0+
 \frac{4s+(5+4K_0)\varepsilon+4K_0h\varepsilon/r_0}{q}.
\end{equation}
\end{lemma}
\begin{proof}
For the measure $\nu$ of $F$, write
\[
 w_a(t)=(1-(t-a)/h)\ind_{(a,a+h)}(t),\qquad
 H(a)=\int w_a(t)\nu(dt)\ge0.
\]
For $0<d<h$, define the nonnegative trapezoid
\[
 T_d(t)=\frac1h\int_{h/2-d}^{h/2-d/2}
                  \ind_{\{|x-t|<h/2\}}\,dx.
\]
Its integral against $\nu$ equals $J(h/2-d/2)-J(h/2-d)$.
Splitting at $-d,-d/2,0,h-d,h-d/2,h$ gives the pointwise bound
\[
 \min\{1,2(h-d)/d\}\,T_d(t)\le w_0(t)+w_{-d}(t).
\]
Integrating and applying the affine approximation to the two $J$ values gives
\begin{equation}\label{eq:two-shifts}
       q\min(d/2,h-d)\le H(0)+H(-d)+2\varepsilon.
\end{equation}
The pointwise bound also holds at the endpoints, with the open intervals
in $w_a$, so atoms require no continuity assumption on $F$.

To transfer over several periods, let
$\overline H_{r_0}(a)=r_0^{-1}\int_a^{a+r_0}H(t)dt$.
Integrating the definition of the uniform convolution gives
\begin{align*}
 \overline H_{r_0}(a+h)-\overline H_{r_0}(a)
 ={}&r_0^{-1}\int_a^{a+r_0}[J(t+3h/2)-J(t+h/2)]dt\\
 &-(h/r_0)[J(a+r_0+h/2)-J(a+h/2)].
\end{align*}
The affine terms cancel, so its magnitude is at most
$2\varepsilon(1+h/r_0)$.
Moreover
$H(a+t)\le H(a)+(t/h)\nu((a+t,a+t+h))$ and the mass is at most
$2(qh+\varepsilon)$, by the affine bound on $J$. Hence
$\overline H_{r_0}(a)\le H(a)+qr_0+\varepsilon r_0/h$.
Write $d=kh+d_0$, $0\le d_0<h$. Transfer from $-d$ to $-d_0$
using the preceding difference bound. If $d_0$ is within $2r_0$
of an endpoint, the conclusion is immediate. Otherwise choose a point
in the averaging interval where $H$ is at most its average and apply
\eqref{eq:two-shifts}, using $H(0)\le s$.
The threshold displacement is at most $r_0$; the resulting bound is
\eqref{eq:spacing-bound}. All arguments of $J$ stay in the stated window.
\end{proof}

\subsection{The scale and the maximizing threshold}
We prove Proposition~\ref{prop:reduction}. Retain the notation of
Section~\ref{sec:selection} for the actual selected maximizer.
For scaled coordinates write
\[
 Y_i=X_i/\ell,\quad v_i=\E Y_i^2,\quad M_i=\E[Y_i|Y_i|],
 \quad\theta=\ell^{-1}\sum_i\E X_i^3,\quad |\theta|\le1.
\]
Let $\rho,r,T,\eps,h$ be as in Proposition~\ref{prop:jitter}.
We give the error quantities explicitly so their later specialization
requires no unspecified constants:
\begin{equation}\label{eq:confinement-errors}
\begin{gathered}
\begin{aligned}
 u&=\rho+\ell/2,& D&=16u+256r,\\
 H_h&=72r+5b_D,& Z^2&=10(u+18r),\\
 S_0&=2\rho+18r+460\ell^2,&
 A_0&=2\rho+24\{4\ell^2+(Z+32\ell)^2/2+\ell/5\},\\
 \gamma&=300\sqrt{A_0}+11S_0+131A_0,&
 W_0&=21D+3\sqrt D+42r,
\end{aligned}\\
\begin{aligned}
 E_P&=14000(u+18r+\eta)+1100Z^2+4000\ell+4000\ell^2+1100\eta,\\
 E_{\rm end}&=E_P+6100\gamma+700\gamma^2+500H_h.
\end{aligned}
\end{gathered}
\end{equation}
Here $a_D,b_D,k_D$ are those in \eqref{eq:lattice-ratios}.
The proof below uses only
\begin{equation}\label{eq:confinement-gates}
 \begin{gathered}
 r\le10^{-6},\quad D\le10^{-3},\quad A_0\le9/64,\quad
 \gamma\le1/100,\\
 E_P\le10^{-4},\quad E_{\rm end}\le1/2000,\quad W_0\le1/4.
 \end{gathered}
\end{equation}
These will follow from \eqref{eq:reduction-parameters} at the end.

The case $h=0$ would give $R\le\phz/6+\rho<\be$ by
\eqref{eq:jitter}, so $h>0$. Initially $h\le128\pi<404$ and
$g(2\pi/h)\le r^2/4$. Equations
\eqref{eq:rounding}--\eqref{eq:rounding-moments} give
$\sqrt e\le101r$, $|Q_i|\le412$.
If $V_Q=\ell^2\sum\E Q_i^2$ and
$B_Q=\ell^2\sum\E|Q_i|^3$, then
$V_Q\ge1-202r$ and $B_Q\le1+125000r$.
Since $B_Q/(hV_Q)\ge1/2$, it follows that
\[
 h\le2(1+125000r)/(1-202r)<5/2.
\]
Apply the same rounding bounds with this improved span. Writing
$K_Q=\ell^2\sum\E Q_i^3$, we obtain
\begin{equation}\label{eq:rounded-data}
 \sqrt e\le r,\quad|V_Q-1|\le3r,\quad
 |B_Q-1|\le32r,\quad|K_Q-\theta|\le32r.
\end{equation}
The lattice inequality in either orientation implies
\begin{equation}\label{eq:span-envelope}
                3h+|\theta|\le c+256r.
\end{equation}
The moment-error allowance is $[32(c+1)+9h]r<256r$.

Unsmoothing at the actual threshold $z$, with shift $h\ell/2$, gives
\begin{equation}\label{eq:global-envelope}
 R\le\phi(z)\{h/2+\theta(1-z^2)/6\}+u.
\end{equation}
The Gaussian increment has error at most $\phz h^2\ell^2/8$,
and the skew correction adds at most $h\ell^2/18$;
together they are below $\ell^2/2$.
Since $\phi(z)$ and $|(1-z^2)\phi(z)|$ are at most $\phz$,
\begin{equation}\label{eq:near-max-envelope}
 3h+|\theta|>c-16u,\qquad R\le\be+u+18r.
\end{equation}
It follows that $h>3/2$. Orient the rounded variables by
$\operatorname{sign}\theta$, arbitrarily if $\theta=0$.
Equations \eqref{eq:rounded-data}--\eqref{eq:near-max-envelope}
bound their aggregate lattice deficit by $D$, its denominator
$hV_Q$ being greater than one. The moment-ratio estimates yield
\begin{equation}\label{eq:near-esseen}
 |h-\he|\le H_h,\qquad
 \bigl||\theta|-1/\sqrt{10}\bigr|\le34r+3k_D+H_h/6.
\end{equation}
For $D\le10^{-3}$,
$a_D\le\sqrt D/3$,
$b_D\le D/2+\sqrt D/9$, and
$k_D\le D/18+2\sqrt D/3$.
Thus $H_h<1/40$ and the second error is below $1/10$.
Since $31/16<\he<41/21$, this proves
\begin{equation}\label{eq:span-range}
                  19/10<h<2,\qquad|\theta|>1/5.
\end{equation}
If $\theta<0$, the envelope in \eqref{eq:global-envelope} has its
maximum at zero because $h>|\theta|$; it would imply
$|\theta|<128r+8u=D/2$, a contradiction. Hence $\theta>0$.
Finally
$\be<(\be+18r)e^{-z^2/2}+u$, so
\begin{equation}\label{eq:threshold-bound}
           |z|\le Z,\qquad0\le K-\be\le u+18r+\eta.
\end{equation}
For this last bound use $-\log(1-x)\le2x$ for $x\le1/2$.

\subsection{Deletion with the same period and the support spacing}
Deleting coordinate $i$ subtracts at most $\ell^2$ from $g$.
Thus its high-frequency product bound is at most $e^{1/2}<2$ times
the full-row bound. Keep the original $h$ for every deletion.
Write $\sigma_i^2=1-s_i\ge3/4$, $\kappa_- =\sum_{j\ne i}\E X_j^3$,
let $J_i$ be the CDF of that deletion plus the same jitter, and put
\[
 G_i^*(x)=\Phi(x/\sigma_i)
       +\frac{\kappa_-}{6\sigma_i^3}(1-(x/\sigma_i)^2)\phi(x/\sigma_i).
\]
The logarithmic remainder in Appendix~\ref{app:fourier} remains
$(49/3)\ell^2t^4$. The low-frequency error is bounded by
$\ell^2e^{-t^2/5}\{(49/3)t^4+t^6/72\}$, whose positive integral
is below $206\ell^2$. The jitter-factor error is below $\ell^2$.
For $a=\delta/\ell$, the Gaussian comparison tail is at most
$4/(3a^2)+\ell/(2a)<5500\ell^2$; use
$\int_a^\infty t^2e^{-3t^2/8}dt\le3/a$.
Also $G_i^*$ is Lipschitz with constant below $1/2$.
The weaker signed smoothing bound therefore gives
\begin{equation}\label{eq:deleted-jitter}
 \|J_i-G_i^*\|_\infty
 \le6000\ell^2+24r\ell/\delta^2
       +2\log(T/\delta)e^{-\eps^2/(2\ell^2)}+20\ell/T
 \le2\rho\ell.
\end{equation}
The last inequality uses the four allocations in
\eqref{eq:jitter-parameters}.

At any original support point $y=\ell t$, the contact equality and
\eqref{eq:gaussian-global} give
\begin{equation}\label{eq:contact-flat}
 |G_i(z-y)-\Phi((z-y)/\sigma_i)-R\ell|\le450\ell^3.
\end{equation}
For the constant, $|t|\le8$, $v_i,|M_i|\le9$,
$\E|Y_i|^3\le27$, and $K\le23/50$.
The Gaussian remainder contributes less than $51\ell^3$,
the quadratic $K$ term less than $51\ell^3$, and the remaining
cubic, moment, and linear terms less than $348\ell^3$.
Equations \eqref{eq:span-envelope}, \eqref{eq:deleted-jitter}, and
\eqref{eq:contact-flat} then imply
\begin{equation}\label{eq:smoothing-loss}
 0\le[J_i(z-y+h\ell/2)-G_i(z-y)]/\ell\le S_0.
\end{equation}
The scale corrections cost at most $10\ell^2$ after division by
$\ell$: use $\sigma_i^{-1}-1\le s_i$,
$\sigma_i^{-3}-1\le3s_i$, and
$|\kappa_-/\ell|\le\theta+27\ell^2$.

Fix a support point $x$. The function
$\ell^{-1}J_i(z-x+\ell a)$ differs from an affine function of slope
$\phz$ by at most $A_0$ on $[-16-2h,3h]$.
Indeed $|a|\le24$, $|z-x+\ell a|\le Z+32\ell$, and
\[
                 |(G_i^*)'(w)-\phz|\le4\ell^2+w^2/2+\ell/5.
\]
Integrate and use \eqref{eq:deleted-jitter}.
Apply Lemma~\ref{lem:support-spacing} to the scaled positive measure,
using \eqref{eq:smoothing-loss}, $q=\phz$, $K_0\le11$, and
$r_0=\sqrt{A_0}\le h/4$. Its bound is at most
$3\sqrt{A_0}+(8/3)(4S_0+49A_0+110\sqrt{A_0})<\gamma$.
Thus for every pair of support points of each original coordinate,
\begin{equation}\label{eq:arithmetic-spacing}
                   \dist((y-x)/\ell,h\Z)\le\gamma.
\end{equation}

\subsection{The support polynomial and the two classes}
Define
\[
 P_i(t)=t^3-c|t|^3+(3c/2)t^2+3(cM_i-v_i)t.
\]
Every point of the full scaled support satisfies
\begin{equation}\label{eq:support-polynomial}
                             P_i(t)\ge-E_P.
\end{equation}
To verify this with constants, divide \eqref{eq:deletion} by
$\phz\ell^3/6$ and use \eqref{eq:gaussian-local}.
Its Taylor remainder divided by $\ell^3$ is below $224\ell$.
Furthermore $|\phi''(z)+\phz|\le(3/2)\phz Z^2$,
since $\|\phi''''\|_\infty\le3\phz$ and $\phi'''(0)=0$.
The respective error bounds are
\[
 13184(K-\be),\quad1092Z^2,\quad3584\ell,\quad
                         3500\ell^2,\quad1008\eta.
\]
For the first, use
$|-|t|^3+(3/2)t^2+3M_it|\le824$ and $6/\phz<16$.
The last two come from $6R\ell s_i^2+\eta b_i+4\eta\ell s_i$.
Equation~\eqref{eq:threshold-bound} makes their sum at most $E_P$.

For a nondegenerate coordinate let $-a,b$ be its support extremes.
Centering gives $v_i\le ab$ and $|M_i|\le v_i$.
If $a,b\ge1/100$, divide the endpoint inequalities by $a,b$ and
add. The linear terms cancel, giving
\[
 (c+1)a^2+(c-1)b^2\le(3c/2)(a+b)+E_P(a+b)/(ab).
\]
The left side is at least $3(a+b)^2$, because
$(c^2-1)/(2c)=3$. Therefore
$a+b\le c/2+E_P/(3ab)\le25/8+1/3<7/2$.
If $a<1/100$ and $a+b\ge7/2$, then $b>3.49$ and
\[
 P_i(b)\le b^2[-(c-1)b+3c/2+3(c-1)a]<-1,
\]
contradicting \eqref{eq:support-polynomial}. If $b<1/100$, use
$P_i(-a)\le a^2[-(c+1)a+3c/2+3(c+1)b]<-1$.
Every support diameter is consequently below $7/2$.

Partition coordinates according as their scaled support diameter is
less than $h/2$ or at least $h/2$. In the first class,
\eqref{eq:arithmetic-spacing} implies diameter at most $\gamma$,
so centering gives $|Y_i|\le\gamma$.
In the second class the diameter differs from $h$ by at most $\gamma$:
zero and two multiples are excluded by \eqref{eq:span-range} and
the diameter bound. Every support point lies within $2\gamma$ of
an extreme. Split at the midpoint, use its upper-cluster indicator
$B_i$, and take conditional means. This gives exactly
\begin{equation}\label{eq:cluster-representation}
 Y_i=h_i(B_i-p_i)+U_i,\quad\E(U_i\mid B_i)=0,\quad
 |h_i-h|\le5\gamma,\quad |U_i|\le2\gamma.
\end{equation}
Both cluster probabilities are positive, because both extremes are
support points and the clusters are separated.

We next bound those probabilities away from zero and one.
On $[-8,8]$, $|P_i'|<1800$. Moving an extreme to its cluster mean
costs at most $3600\gamma$. Conditional centering gives
$|M_i-M_i^0|\le\E U_i^2\le4\gamma^2$ and
$v_i-v_i^0=\E U_i^2\le4\gamma^2$, where the superscript denotes
the pure two-point law. The first bound uses the Lipschitz derivative
of $x|x|$ and also holds if a cluster crosses zero.
Their polynomial contribution is below $700\gamma^2$.
At a pure-law endpoint the polynomial is
$h_i^3C_3(p_i)+h_i^2C_2(p_i)$, with $|C_3|<29$, $|C_2|<75/8$.
Its gap derivative for gaps at most $9/4$ is below $500$.
Thus moving to the ideal gap $\he$ leaves both endpoint values at
least $-E_{\rm end}$.

Put $t_0=3\sqrt{10}/2$. At its upper and lower endpoints those
ideal-gap values are
\[
 \he^3(1-p)^2A(p),\qquad\he^3p^2B(p),
\]
where
\begin{align*}
 A(p)&=t_0-(c-1)+4(c-1)p-6cp^2,\\
 B(p)&=t_0-3(c-1)+(8c-4)p-6cp^2.
\end{align*}
For $p\le1/4$, $B$ is increasing and
$B(p)\le(\sqrt{10}-17)/8<-1$.
If $p\in[1/100,1/4]$, the lower endpoint is below
$-6/10000<-1/2000$, since $\he^3>6$.
For $p\le1/100$, use $A(0)=-\pe<-2/5$, $A'\le21$, so
$A(p)<-19/100$. For $p\ge3/4$, $A$ is decreasing and
$A(p)\le(\sqrt{10}-49)/8<-5$, excluding $[3/4,99/100]$.
For $p\ge99/100$, $B(1)<-2$ and $|B'|\le46$ give $B(p)<-1$.
In view of $E_{\rm end}\le1/2000$, all these cases are impossible.
Hence every two-cluster probability lies in $[1/4,3/4]$.

\subsection{The extra coordinates and the parameter check}
Lattice stability and the rounding coupling give
\begin{equation}\label{eq:original-W1}
 \sum_i\ell^2v_i W_1(\Law(Y_i),\Law(h(B_{\pe}-\pe)))\le W_0.
\end{equation}
For the rounded variables \eqref{eq:lattice-W1} is at most
$(1+3r)(20D+2\sqrt D)\le21D+3\sqrt D$.
The change of variance weights has total absolute mass at most $3r$,
and the distances are at most $13$, costing $39r$.
The original variance-weighted rounding distance is at most $3r$,
by $v_i\le9$ and the squared error bound $e\le r^2$.
This proves \eqref{eq:original-W1} with the $42r$ allowance in $W_0$.

The target law has absolute first moment
$2h\pe\qe>9/10$. On the small-diameter class $|Y_i|\le\gamma$.
Its total original variance $V_s$ is therefore at most
$W_0/(9/10-\gamma)\le2W_0$.
Each remaining coordinate has scaled variance between $1/2$ and $2$,
by \eqref{eq:span-range}, \eqref{eq:cluster-representation}, and
$p_i\in[1/4,3/4]$. Their number $m$ satisfies
\begin{equation}\label{eq:core-count}
 \frac{1-2W_0}{2\ell^2}\le m\le\frac2{\ell^2},
                         \qquad m\ge1/(4\ell^2).
\end{equation}
Calling the small scaled coordinates $Z_j$,
\[
 \frac{\sum_j\Var Z_j}{mh^2}
       \le\frac{4W_0}{h^2(1-2W_0)}<3W_0.
\]
The relative gap errors in \eqref{eq:cluster-representation} are
below $3\gamma$, the residuals satisfy $|U_i|/h<2\gamma$, and
$|Z_j|/h<\gamma$.
Thus the conclusion follows once $3\gamma,3W_0\le\lambda$ and
$\ell\le1/(2\sqrt N)$.

We finish by checking these conditions and
\eqref{eq:confinement-gates}. From \eqref{eq:jitter-parameters},
$\ell\le\rho/48000$ and $r=\rho/196608$. For $\rho\le10^{-6}$,
\begin{gather*}
 D\le17\rho,\quad Z^2\le20\rho,\quad
 H_h\le43\rho+3\sqrt\rho,\quad S_0\le3\rho,\quad A_0\le500\rho,\\
 \gamma\le7000\sqrt\rho,\quad W_0\le14\sqrt\rho,
             \quad E_P\le51000\rho+15100\eta.
\end{gather*}
For $\rho\le10^{-12}$ these imply
\[
 E_{\rm end}\le5\cdot10^7\sqrt\rho+2\cdot10^4\eta.
\]
Before enlargement its $\rho$ and $\sqrt\rho$ coefficients are
at most $34300072500$ and $42701500$, which proves this last bound.
Substitute $\rho=10^{-26}\lambda^2$, $\eta=10^{-10}\lambda$:
\[
 E_{\rm end}\le7\cdot10^{-6}\lambda<1/2000,\quad
 \gamma\le7\cdot10^{-10}\lambda,\quad
 W_0\le1.4\cdot10^{-12}\lambda.
\]
All the stated conditions follow, including $E_P\le10^{-4}$.
The representations and the variance estimates concern the original
coordinates throughout, proving Proposition~\ref{prop:reduction}.

\section{Conditional moments and the one-sided comparison}
\label{app:conditional}
This appendix proves the estimates used in
Section~\ref{sec:local}. All conditioning is on the actual Bernoulli
count. Dependence between a label and its residual is preserved.

\subsection{Exponential tilting and projection}
Retain $p_i\in[1/4,3/4]$, $v_i=p_i(1-p_i)$,
$\sum_i v_i a_i=0$, and $\Lambda_A=\sum_i v_i a_i^2$.
Suppose $m\ge10^{12}$ and $|k-\mu|\le6\sqrt m$.
Tilt by $e^{\theta K_0}$ so that the count mean is $k$, putting
\[
 r_i=\frac{p_i e^\theta}{1-p_i+p_i e^\theta},\quad
 w_i=r_i(1-r_i),\quad V_\theta=\sum_i w_i.
\]
For $|\theta|\le\log2$, $r_i\in[1/8,7/8]$ and $w_i\ge7/64>1/10$.
The count mean changes by more than $6\sqrt m$ in both directions
before that boundary, so the tilt exists and
\begin{equation}\label{eq:tilt-quantities}
 |\theta|\le60/\sqrt m,\quad |w_i/v_i-1|\le2|\theta|,
       \quad |r_i-p_i-\theta v_i|\le\theta^2v_i.
\end{equation}
The conditional law given $K_0=k$ is unchanged. Define
\[
 c_\theta=\frac{\sum_iw_i a_i}{V_\theta},\quad b_i=a_i-c_\theta,
 \quad L=\sum_i b_i(B_i-r_i),\quad Q=\sum_i b_i^2,
 \quad Q_\theta=\sum_iw_i b_i^2.
\]
On the count event, $A=m_k^0+L$, where
\begin{equation}\label{eq:tilted-gap-data}
 |m_k^0|=\left|\sum_i a_i(r_i-p_i)\right|
      \le1800\sqrt{\Lambda_A}/\sqrt m,\quad
 |Q_\theta-\Lambda_A|\le180\Lambda_A/\sqrt m,
 \quad Q\le16\Lambda_A.
\end{equation}
For the second bound expand
$Q_\theta=\sum_iw_i a_i^2-(\sum_iw_i a_i)^2/V_\theta$.
The errors are at most $2|\theta|\Lambda_A$ and
$8\theta^2\Lambda_A$, together at most $3|\theta|\Lambda_A$.
Also $\Lambda_A/2\le Q_\theta\le3\Lambda_A/2$; using
$w_i\ge7/64$ proves the last bound. If $\Lambda_A=0$, all the
gap estimates are immediate, so suppose it is positive.

\subsection{Fourier moments with their residual factors}
For $|t|\le1/4$, put
$z_i(t)=r_ie^{it}/(1-r_i+r_ie^{it})$.
The first centered complex cumulant is $c_{i,1}=z_i-r_i$; the next
three are
\[
 c_{i,2}=z_i(1-z_i),\quad
 c_{i,3}=z_i(1-z_i)(1-2z_i),\quad
 c_{i,4}=z_i(1-z_i)(1-6z_i(1-z_i)).
\]
Both $z_i$ and $1-z_i$ have modulus at most two. Differentiation gives
\[
 |c_{i,1}-iw_it|\le10t^2,\quad |c_{i,2}-w_i|\le20|t|,
 \quad |c_{i,3}|\le20,\quad |c_{i,4}|\le100.
\]
For $M_j=\sum_i b_i^jc_{i,j}$, the exact cancellation
$\sum_iw_i b_i=0$ yields
\begin{gather}
 |M_1|\le10\sqrt{mQ}t^2,\quad
 |M_2-Q_\theta|\le20Q|t|,\quad |M_2|\le6Q,\notag\\
 |M_3|\le20Q^{3/2},\qquad |M_4|\le100Q^2.              \label{eq:complex-moments}
\end{gather}
The tilted count transform $F_\theta$ satisfies
$|F_\theta(t)|\le e^{-mt^2/50}$ on $[-\pi,\pi]$.
Lemma~\ref{lem:count-estimate} gives
$\Pp_\theta(K_0=k)\ge1/(3\sqrt m)$.
After division by this probability, upper bounds for the Gaussian
integrals of powers $0,1,2,4,8$ are respectively
\begin{equation}\label{eq:gaussian-moment-table}
             7,\quad25m^{-1/2},\quad175m^{-1},\quad
                          13125m^{-2},\quad3\cdot10^8m^{-4}.
\end{equation}
These follow by scaling $t\sqrt m$ in
$(2\pi)^{-1}\int e^{-mt^2/50}|t|^jdt$.
On $1/4\le|t|\le\pi$, expand the moment before factoring the
product. A term of degree $j\le4$ removes at most four Bernoulli
factors; the remaining product is at most $e^{-m/810}$.
The conditional tail contribution is therefore at most
$3m^{(j+1)/2}Q^{j/2}e^{-m/810}$.
For $j=1,2,4$ and $m\ge10^8$, these are at most
$\sqrt Q/\sqrt m$, $Q/\sqrt m$, and $Q^2$.
For instance $e^{-m/810}\le10^{18}/m^5$ reduces the verification
to decreasing powers of $m$.

The first, second, and fourth moment numerators are exactly
\[
 F_\theta M_1,\quad F_\theta(M_1^2+M_2),\quad
 F_\theta(M_1^4+6M_1^2M_2+3M_2^2+4M_1M_3+M_4).
\]
Equations \eqref{eq:complex-moments}--\eqref{eq:gaussian-moment-table},
with the tails just bounded, imply
\begin{equation}\label{eq:projected-moments}
 \begin{gathered}
 |\E(L\mid K_0=k)|\le2000\sqrt Q/\sqrt m,\\
 |\Var(L\mid K_0=k)-Q_\theta|\le2000Q/\sqrt m,
 \qquad \E(L^4\mid K_0=k)\le2000Q^2.
 \end{gathered}
\end{equation}
Here is numerical detail for these enlargements.
The mean's central coefficient is $1750$.
The second moment errors are at most
$500Q/\sqrt m+1312500Q/m+2Q/\sqrt m$,
and the squared mean costs at most $4\cdot10^6Q/m$.
The five fourth-moment central coefficients, in the displayed order,
are at most
\[
 3\cdot10^{12}/m^2,\quad47250000/m,\quad756,\quad
                         140000/\sqrt m,\quad700,
\]
all multiplying $Q^2$. Their sum and the tail are below $2000Q^2$.
Centering raises the fourth-moment bound to $20000Q^2$.
Combining \eqref{eq:tilted-gap-data} and
\eqref{eq:projected-moments} proves
\eqref{eq:conditional-mean}--\eqref{eq:conditional-gap}.

For completeness, the relative label estimate
\eqref{eq:conditional-label} also follows numerically.
Under the tilt, deleting a Bernoulli changes count variance by at most
$1/4$ and displaces the new threshold from its mean by at most one.
The Gaussian mass ratio has logarithm bounded by $1/V_\theta$.
The $6/V$ errors and a lower bound by half the centered Gaussian mass
then give a deleted-count ratio error at most
\[
 4/V_\theta+28/(\phz\sqrt{V_\theta})\le300/\sqrt m.
\]
Each tilted label probability differs relatively from its original
probability by at most $2|\theta|\le120/\sqrt m$.
Multiplication gives $500/\sqrt m$ for $m\ge10^{12}$.
This is a relative bound, so arbitrarily small within-cluster noise
variances remain controlled after conditioning.

\subsection{Proof of the one-sided loss}
We prove \eqref{eq:one-sided}, following the clipping method of
\citet[Lemma~3.2 and equation~(83)]{HeCheng2026}.
First suppose $V_0$ is centered and $|V_0|\le a=1/10$.
For $u\ge0$, $\E(V_0)_+\le u+a\Pp(V_0>u)$ and
$f(u)\ge3u/4$ give
$\Pp(V_0>u)+f(u)\ge3\E|V_0|/8$.
If $u=-r$, $0<r\le a$, put $t=\Pp(V_0>-r)$.
Centering gives $0\le at-r(1-t)$, hence $r\le2at$.
Also $\E|V_0|=2\E(V_0)_+\le2at$, so
$t+f(-r)\ge t-5r/4\ge3t/4\ge3\E|V_0|/8$.
For $a<r\le13/20$, the probability is one and
$1+f(-r)\ge3/16\ge3\E|V_0|/8$.

For general centered $Y$, clip at $1/20$, obtaining $T$, and put
$b=\E T$, $V_0=T-b$, $M_4=\E Y^4$.
Then
\[
 |b|\le8000M_4,\quad\Pp(T\ne Y)\le160000M_4,\quad
                  \E|V_0|\ge\E|Y|-16000M_4.
\]
If $M_4\le1/160000$, then $|b|\le1/20$,
$|V_0|\le1/10$, and $|u-b|\le13/20$.
Apply the bounded case at $u-b$, and compare $f(u-b)$ with $f(u)$.
The total error coefficient is
$(3/8)16000+160000+(5/4)8000=176000$.
This proves the first inequality in \eqref{eq:one-sided}.
Apply it to $-Y$ and $-f(-\,\cdot\,)$ for the second.
All events in this proof are strict; atoms at the threshold are included
with the conventions needed in \eqref{eq:layers}.

\bibliographystyle{plainnat}
\bibliography{refs}
\end{document}
